\documentclass[11pt,a4paper]{report}

\usepackage[T1]{fontenc}
\usepackage[utf8]{inputenc}
\usepackage{lmodern}
\usepackage{amsmath}
\usepackage{amssymb}
\usepackage[margin=2.5cm]{geometry}
\usepackage{booktabs}
\usepackage{longtable}
\usepackage{array}
\usepackage{xcolor}
\usepackage{listings}
\usepackage{fancyvrb}
\usepackage{enumitem}
\usepackage{titlesec}
\usepackage[hidelinks,breaklinks]{hyperref}
\usepackage{tcolorbox}
\tcbuselibrary{skins,breakable}

\definecolor{codebg}{HTML}{F6F8FA}
\definecolor{codeframe}{HTML}{D8DEE4}
\definecolor{accent}{HTML}{1F4E79}
\definecolor{warnbg}{HTML}{FFF6E5}
\definecolor{warnframe}{HTML}{E0A800}
\definecolor{okbg}{HTML}{EEF7EE}
\definecolor{okframe}{HTML}{4C8C4A}

\lstset{
  basicstyle=\ttfamily\small,
  backgroundcolor=\color{codebg},
  frame=single,
  rulecolor=\color{codeframe},
  framesep=5pt,
  breaklines=true,
  breakatwhitespace=false,
  columns=fullflexible,
  keepspaces=true,
  showstringspaces=false,
  literate={-}{{-}}1 {_}{{\_}}1
}

\newtcolorbox{note}[1][]{colback=okbg,colframe=okframe,boxrule=0.6pt,
  left=6pt,right=6pt,top=5pt,bottom=5pt,breakable,#1}
\newtcolorbox{warn}[1][]{colback=warnbg,colframe=warnframe,boxrule=0.6pt,
  left=6pt,right=6pt,top=5pt,bottom=5pt,breakable,#1}

\titleformat{\chapter}[display]
  {\normalfont\huge\bfseries\color{accent}}{\chaptertitlename\ \thechapter}{16pt}{\Huge}
\setlist{itemsep=2pt,parsep=0pt,topsep=4pt}

\newcommand{\econenvversion}{1.0.8}

\title{
  \vspace{-1.5cm}
  {\Huge\bfseries EconEnv}\\[8pt]
  {\Large One Notebook. Multiple Econometric Engines.}\\[18pt]
  {\large Python, R, Stata and EViews in a single Jupyter workflow}\\[26pt]
  {\large\bfseries Installation and User Guide}\\[10pt]
  {\normalsize for EconEnv \econenvversion}
}
\author{
  {\large Developed by}\\[6pt]
  {\Large\bfseries Dr Merwan Roudane}\\[8pt]
  {\normalsize \url{https://github.com/merwanroudane}}
}
\date{\today}

\begin{document}
\maketitle

\begin{center}
\begin{minipage}{0.86\textwidth}
\small
\textbf{About this guide.} It assumes nothing. Chapter~\ref{ch:install} starts
from a computer with no Python on it and finishes with all four programs
talking to each other in one notebook. If you already have Python and Jupyter,
skip to Section~\ref{sec:installecon}.

\medskip
\textbf{EconEnv does not contain, bundle or redistribute Stata or EViews.}
Those are commercial products which you must install and license yourself. R and
Python are free and open source. EconEnv only connects to what is already on
your machine.
\end{minipage}
\end{center}

\vspace{10pt}
\begin{center}
\begin{minipage}{0.86\textwidth}
\textbf{Quick start.} If you already have Python:
\begin{lstlisting}
pip install econenv
\end{lstlisting}
then, in a notebook:
\begin{lstlisting}
%load_ext econenv
%econ doctor
\end{lstlisting}
\end{minipage}
\end{center}

\vfill
\begin{center}
\small
Website: \url{https://merwanroudane.github.io/econenv/}\\[3pt]
Repository: \url{https://github.com/merwanroudane/econenv} \quad$\cdot$\quad
Package: \url{https://pypi.org/project/econenv/\econenvversion/}
\end{center}

\tableofcontents

% ======================================================================
\chapter{What EconEnv is, and why}

\section{The problem}

An applied econometrics paper rarely lives inside one program. The unit-root
test is done in EViews because that is where the output is readable. The panel
estimator is in Stata because \texttt{xtreg} is the reference implementation.
The plots are in R because \texttt{ggplot2} is better. The data cleaning is in
Python because pandas is better.

So the working day looks like this:

\begin{center}
\texttt{Python $\rightarrow$ to\_csv() $\rightarrow$ Stata $\rightarrow$ export
delimited $\rightarrow$ R $\rightarrow$ write.csv $\rightarrow$ EViews}
\end{center}

Four programs open, four windows, four copies of the same data slowly drifting
apart. A missing value that meant \texttt{.a} in Stata arrives as an empty cell
in R. A quarterly index becomes a string. And when a referee asks why your
robust standard error differs from theirs, there is no way to answer without
redoing the whole chain by hand.

\section{What EconEnv does}

EconEnv turns the four programs into \emph{execution engines behind one Python
kernel}. You write:

\begin{lstlisting}
%load_ext econenv

df = pd.read_csv("data.csv")        # Python, as usual
\end{lstlisting}

\begin{lstlisting}
%%R -i df
fit <- lm(y ~ x1 + x2, data = df)
summary(fit)
\end{lstlisting}

\begin{lstlisting}
%%stata
regress y x1 x2, robust
\end{lstlisting}

\begin{lstlisting}
%%eviews -i df
equation eq1.ls y c x1 x2
eq1.output
\end{lstlisting}

One notebook. One kernel. One dataset. No CSV round-trip.

And then the part that is genuinely hard to do any other way:

\begin{lstlisting}
econenv.compare_ols(df, "y ~ x1 + x2")
\end{lstlisting}

which estimates the same specification in every engine you have and puts the
results side by side --- reporting where they differ, and why, instead of
quietly picking one.

\section{What it is not}

\begin{itemize}
  \item It is \textbf{not} a new Jupyter kernel. It is a Python package and an
        IPython extension. Your notebook stays a normal Python notebook.
  \item It does \textbf{not} reimplement \texttt{\%\%stata}. Stata~17 and later
        ship official IPython magics, and EconEnv loads those rather than
        writing worse ones.
  \item It does \textbf{not} ship any commercial software, licence file or
        binary.
  \item It does \textbf{not} try to make the programs agree. Where their
        defaults differ, it says so.
\end{itemize}

% ======================================================================
\chapter{Requirements}

\section{Minimum}

\begin{center}
\begin{tabular}{@{}lll@{}}
\toprule
\textbf{Component} & \textbf{Version} & \textbf{Needed for} \\
\midrule
Python & 3.9 or newer & everything \\
pandas, numpy, IPython & installed automatically & everything \\
statsmodels & installed automatically & the Python engine \\
JupyterLab or Notebook & any recent & running notebooks \\
\bottomrule
\end{tabular}
\end{center}

\section{Optional engines}

Install only what you use. None of these is required to install EconEnv.

\begin{center}
\begin{tabular}{@{}llp{6.6cm}@{}}
\toprule
\textbf{Engine} & \textbf{Version} & \textbf{Notes} \\
\midrule
R & 4.0+ & Free. EconEnv finds it automatically; no \texttt{PATH} setup needed. \\
Stata & \textbf{17 or newer} & Commercial. PyStata ships with Stata~17+; earlier versions cannot be driven this way. \\
EViews & 12, 13 or 14 & Commercial. \textbf{Windows only} --- automation is COM-based. \\
\bottomrule
\end{tabular}
\end{center}

\section{Operating systems}

\begin{center}
\begin{tabular}{@{}llll@{}}
\toprule
 & \textbf{Windows} & \textbf{Linux} & \textbf{macOS} \\
\midrule
Python & yes & yes & yes \\
R      & yes & yes & yes \\
Stata  & yes & yes & yes \\
EViews & yes & no  & no \\
\bottomrule
\end{tabular}
\end{center}

\begin{warn}
EViews automation uses Microsoft COM, which exists only on Windows. On Linux
and macOS EconEnv installs and runs normally; the EViews engine simply reports
itself unavailable and everything else works. Linux and macOS are supported by
design but have not yet been exercised by the author --- reports are welcome.
\end{warn}

% ======================================================================
\chapter{Installation from zero}
\label{ch:install}

This chapter assumes a Windows machine with nothing installed. Adapt the paths
if yours differ.

\section{Step 1 --- Python}

Two reasonable routes.

\subsection*{Route A: python.org (simplest)}

\begin{enumerate}
  \item Go to \url{https://www.python.org/downloads/} and download Python 3.11
        or 3.12 for Windows.
  \item Run the installer. \textbf{Tick ``Add python.exe to PATH''} on the first
        screen --- this is the single most common thing people miss.
  \item Open a new Command Prompt and check:
\end{enumerate}

\begin{lstlisting}
python --version
\end{lstlisting}

\subsection*{Route B: Anaconda / Miniconda}

Convenient if you also want conda packages. Download from
\url{https://www.anaconda.com/download}, then create a dedicated environment so
you do not disturb the base one:

\begin{lstlisting}
conda create -n econ python=3.12
conda activate econ
\end{lstlisting}

\begin{note}
Whichever route you choose, remember \emph{which} Python it is. If you install
EconEnv into one interpreter and start Jupyter from another, the extension will
not be found --- and the error will not obviously say why. Section~\ref{sec:wrongpython}
covers this.
\end{note}

\section{Step 2 --- Jupyter}

\begin{lstlisting}
pip install jupyterlab
\end{lstlisting}

Start it with:

\begin{lstlisting}
jupyter lab
\end{lstlisting}

\section{Step 3 --- EconEnv}
\label{sec:installecon}

\begin{lstlisting}
pip install econenv
\end{lstlisting}

That is the whole core installation. It pulls in pandas, numpy, IPython and
statsmodels, and nothing commercial.

Optional extras, installed only if you need them:

\begin{lstlisting}
pip install "econenv[eviews]"   # comtypes, for EViews (Windows)
pip install "econenv[arrow]"    # faster data transfer to R
pip install "econenv[stata]"    # helper for locating PyStata
pip install "econenv[all]"      # all of the above
\end{lstlisting}

Check it worked:

\begin{lstlisting}
econenv doctor
\end{lstlisting}

This prints a diagnosis of every layer --- Python, Jupyter, R, Stata, EViews ---
and, for anything wrong, what to do about it.

\section{Step 4 --- R (optional, free)}

\begin{enumerate}
  \item Download R from \url{https://cran.r-project.org/bin/windows/base/}.
  \item Run the installer and accept the defaults. It lands in
        \texttt{C:\textbackslash Program Files\textbackslash R\textbackslash R-4.x.x}.
  \item Nothing else is needed. EconEnv searches the registry and the usual
        install locations; you do not have to set \texttt{R\_HOME} or touch
        \texttt{PATH}.
\end{enumerate}

If you have several versions installed and want a specific one:

\begin{lstlisting}
%econ config r.home "C:/Program Files/R/R-4.5.2"
\end{lstlisting}

\begin{note}
\textbf{You do not need rpy2.} EconEnv's default R backend is a persistent
\texttt{Rterm} child process, which needs no compiler and no \texttt{R\_HOME}
gymnastics. rpy2 publishes no Windows wheels, so \texttt{pip install rpy2}
typically leaves you with a broken build. If you want it anyway, use
conda-forge --- see Section~\ref{sec:rpy2}.
\end{note}

\section{Step 5 --- Stata (optional, commercial)}

You need \textbf{Stata 17 or newer}, because that is when StataCorp began
shipping PyStata.

\begin{enumerate}
  \item Install and licence Stata normally.
  \item Nothing else. EconEnv finds the installation, validates the executable
        and detects the edition (BE, SE or MP).
\end{enumerate}

To pin a specific installation or edition:

\begin{lstlisting}
%econ config stata.home "C:/Program Files/StataNow19"
%econ config stata.edition mp
\end{lstlisting}

\section{Step 6 --- EViews (optional, commercial, Windows only)}

\begin{enumerate}
  \item Install and licence EViews normally.
  \item \textbf{Open EViews once, by hand, and close it.} This registers its
        COM automation server and validates the licence. Until you have done
        this, automation will not connect.
  \item Install the COM bridge:
\end{enumerate}

\begin{lstlisting}
pip install "econenv[eviews]"
\end{lstlisting}

\begin{warn}
\textbf{If you have more than one EViews version installed}, the generic
\texttt{EViews.Manager} identifier binds to whichever registered \emph{last} ---
not the newest. On a machine with EViews 12, 13 and 14, it may well be 13 that
answers. \texttt{\%econ status} shows the version actually connected. To pin one:

\begin{lstlisting}
%econ config eviews.progid EViews.Manager.14
\end{lstlisting}

Note the format: \texttt{EViews.Manager.14}, not \texttt{EViews14.Manager}.
\end{warn}

\section{Step 7 --- Verify everything}

In a notebook:

\begin{lstlisting}
%load_ext econenv
%econ status
\end{lstlisting}

You should see a table like this:

\begin{center}
\ttfamily\footnotesize
\begin{tabular}{@{}llllll@{}}
\toprule
engine & state & version & edition & backend & location \\
\midrule
Python & configured & 3.11.0 & & cpython & C:\textbackslash ...\textbackslash Python311 \\
EViews & configured & 13 & & comtypes & C:\textbackslash Program Files\textbackslash EViews 13 \\
R & configured & 4.5.2 & & subprocess & C:\textbackslash Program Files\textbackslash R\textbackslash R-4.5.2 \\
Stata & configured & 19.5 & mp & pystata & C:\textbackslash Program Files\textbackslash StataNow19 \\
\bottomrule
\end{tabular}
\end{center}

Then the full diagnosis:

\begin{lstlisting}
%econ doctor
\end{lstlisting}

Every check reports PASS, WARNING or ERROR, and every warning and error carries
a suggested fix.

% ======================================================================
\chapter{First steps}

\section{Loading the extension}

Once per notebook:

\begin{lstlisting}
%load_ext econenv
\end{lstlisting}

It prints which magics it registered and who owns each one:

\begin{lstlisting}
EconEnv 1.0.8 loaded - one notebook, multiple econometric engines.
  %econ                  econenv
  %Rec / %%Rec           econenv
  %R / %%R               econenv (rpy2 not installed)
  %eviews / %%eviews     econenv
  %stata / %%stata       pystata (official)
\end{lstlisting}

\begin{note}
Note the last line. \texttt{\%\%stata} belongs to StataCorp, not to EconEnv.
That means it takes \emph{Stata's} options, not EconEnv's: \texttt{\%\%stata
-{}-result} is a syntax error. Move data with \texttt{\%econ push} instead.
\end{note}

\section{Sending data to an engine}

Three equivalent ways:

\begin{lstlisting}
%econ push r df           # line magic
\end{lstlisting}

\begin{lstlisting}
%%R -i df                 # inline, when running R code anyway
summary(df)
\end{lstlisting}

\begin{lstlisting}
econenv.push("r", "df", df)   # from Python
\end{lstlisting}

\section{Getting results back}

\begin{lstlisting}
back = econenv.pull("stata")          # the current Stata dataset
one  = econenv.pull("eviews", "yf")   # one EViews series
frame = econenv.pull("r", "mydata")   # a named R data frame
\end{lstlisting}

\begin{warn}
R has no notion of a ``current dataset'' the way Stata and EViews do --- every
data frame is just a variable. \texttt{econenv.pull("r")} therefore returns the
frame EconEnv last transferred; if there is none, it tells you so and lists the
frames R actually holds. Name the frame explicitly when in doubt.
\end{warn}

\section{Moving data between two engines directly}

\begin{lstlisting}
econenv.move("stata", "eviews", "mydata")
\end{lstlisting}

No file is written, and no manual step is involved.

% ======================================================================
\chapter{The magic commands}

\section{\texttt{\%econ} --- managing everything}

\begin{center}
\begin{longtable}{@{}lp{9.2cm}@{}}
\toprule
\textbf{Command} & \textbf{What it does} \\
\midrule
\endhead
\texttt{\%econ status} & Table of every engine: state, version, backend, location \\
\texttt{\%econ engines} & The same, in detail, with capabilities \\
\texttt{\%econ doctor} & Full diagnosis with a fix for every problem \\
\texttt{\%econ versions} & Version of EconEnv and of each engine \\
\texttt{\%econ capabilities} & What each engine can do on this machine \\
\texttt{\%econ models} & Which estimators are available in which engine \\
\texttt{\%econ start r} & Start an engine explicitly \\
\texttt{\%econ stop} / \texttt{restart} / \texttt{reset} & Session control \\
\texttt{\%econ config r.home "..."} & Show or set a configuration value \\
\texttt{\%econ push <engine> <name>} & Send a Python object to an engine \\
\texttt{\%econ pull <engine> [name]} & Fetch an object back \\
\texttt{\%econ move <from> <to> <name>} & Engine to engine \\
\texttt{\%econ snapshot} & Reproducibility snapshot of every version \\
\texttt{\%econ eviews [topic]} & Look up EViews commands (see below) \\
\texttt{\%econ ols y \textasciitilde{} x} & Run one OLS across every engine \\
\bottomrule
\end{longtable}
\end{center}

\section{\texttt{\%\%R} --- R}

\begin{lstlisting}
%%R -i df -o results
fit <- lm(y ~ x, data = df)
results <- coef(fit)
\end{lstlisting}

\texttt{-i} sends objects in, \texttt{-o} brings them back. \texttt{-q}
suppresses output, \texttt{-{}-no-graphics} suppresses plot capture.

If rpy2 is installed, EconEnv loads \emph{rpy2's own} \texttt{\%R} magics rather
than shadowing them, and its own are always available as \texttt{\%Rec} /
\texttt{\%\%Rec}.

\section{\texttt{\%\%stata} --- Stata}

This is StataCorp's official magic. Use Stata's own syntax and options:

\begin{lstlisting}
%%stata
regress y x1 x2, robust
estat ic
\end{lstlisting}

\section{\texttt{\%\%eviews} --- EViews}

\begin{lstlisting}
%%eviews -i df
equation eq1.ls y c x1 x2
eq1.output
line y
\end{lstlisting}

\texttt{-i} pushes a DataFrame in as a workfile, \texttt{-o} pulls series out,
\texttt{-p} selects a page, \texttt{-{}-no-graphs} suppresses figures.

% ======================================================================
\chapter{EViews for people who use the GUI}

This is the chapter most EViews users need, because in EViews you normally
click, and in a notebook there is nothing to click.

\section{Look commands up without leaving the notebook}

\begin{lstlisting}
%econ eviews                      # the list of tasks
%econ eviews graph                # everything about plotting
%econ eviews find cointegration   # search all of it
\end{lstlisting}

Each result gives the menu path you already know, the command, what it does, and
an example:

\begin{lstlisting}
g.coint(e)
    Johansen system cointegration test - trace and maximum-eigenvalue
    statistics for how many cointegrating relations exist.
    GUI: Group > View > Cointegration Test > Johansen
    e.g. g1.coint(e)
\end{lstlisting}

The catalogue holds 136 commands, 94 of them verified against EViews 13.

\section{The one rule that makes EViews commands guessable}

EViews objects work like this:

\begin{center}
\texttt{object\_name.what\_you\_want}
\end{center}

\begin{itemize}
  \item Anything in an object's \textbf{View} menu is \texttt{object.thatview}.
  \item Anything in its \textbf{Proc} menu is \texttt{object.thatproc}, usually
        followed by a name for whatever it creates.
\end{itemize}

So \texttt{eq1.output} shows an equation's results, \texttt{x.line} plots a
series, \texttt{v1.impulse} draws impulse responses.

\section{Three ways to draw the same plot}

All three work, and the GUI hides the distinction:

\begin{lstlisting}
line x            % command form  - quickest
x.line            % view form     - matches the View menu
graph gr1.line x  % object form   - keeps the graph so you can edit it
\end{lstlisting}

\section{The commands you will need first}

\begin{center}
\begin{longtable}{@{}llp{5.4cm}@{}}
\toprule
\textbf{You would click} & \textbf{Command} & \textbf{What it does} \\
\midrule
\endhead
File $>$ New $>$ Workfile & \texttt{wfcreate q 1990Q1 2020Q4} & Quarterly workfile \\
Quick $>$ Generate Series & \texttt{series ly = log(y)} & Create or transform \\
Quick $>$ Sample & \texttt{smpl 1995Q1 2015Q4} & Restrict the sample \\
Open as Group & \texttt{group g1 x y} & Bundle series \\
Series $>$ View $>$ Graph $>$ Line & \texttt{x.line} & Plot \\
Descriptive Statistics & \texttt{x.stats} & Summary table \\
Correlogram & \texttt{x.correl} & ACF/PACF with Q-stats \\
Unit Root Test $>$ ADF & \texttt{x.uroot(adf)} & Dickey--Fuller \\
Cointegration $>$ Johansen & \texttt{g1.coint(e)} & Trace and max-eigen \\
Estimate Equation $>$ LS & \texttt{equation eq1.ls y c x} & OLS; \texttt{c} is the constant \\
Estimate Equation $>$ ARCH & \texttt{equation eq1.arch(1,1) y c x} & GARCH(1,1) \\
Estimate VAR & \texttt{var v1.ls 1 2 x y} & VAR, lags 1 to 2 \\
Equation $>$ View $>$ Output & \texttt{eq1.output} & Results table \\
Actual, Fitted, Residual & \texttt{eq1.resids} & The first thing to look at \\
Serial Correlation LM & \texttt{eq1.auto(2)} & Breusch--Godfrey \\
Heteroskedasticity, White & \texttt{eq1.white} & White's test \\
Recursive $>$ CUSUM & \texttt{eq1.rls(q)} & Parameter stability \\
Recursive $>$ CUSUM squares & \texttt{eq1.rls(v)} & Variance stability \\
Forecast & \texttt{eq1.forecast(g) yf} & Forecast and plot \\
\bottomrule
\end{longtable}
\end{center}

The complete reference is in \texttt{docs/engines/eviews-commands.md}.

\begin{note}
\textbf{If a command runs but you see nothing}, that is a bug in EconEnv, not in
what you typed. Every display command should produce output, a plot, or a
warning saying it could not be read. Silence is never correct --- please report
it.
\end{note}

% ======================================================================
\chapter{MATLAB}

MATLAB joins the other four through the official MATLAB Engine API for Python,
which MathWorks ships inside every installation and publishes on PyPI.

\section{Installing the bridge}

The engine package must match the MATLAB release:

\begin{center}
\begin{tabular}{@{}ll@{}}
\toprule
\textbf{MATLAB release} & \textbf{Engine package} \\
\midrule
R2024a & \texttt{pip install "matlabengine==24.1.*"} \\
R2024b & \texttt{pip install "matlabengine==24.2.*"} \\
R2025a & \texttt{pip install "matlabengine==25.1.*"} \\
\bottomrule
\end{tabular}
\end{center}

\begin{warn}
It must also match your \emph{Python}. R2024a's engine supports Python 3.9 to
3.11 and refuses anything newer, so a mismatch is an import error rather than a
wrong answer. \texttt{econenv doctor} says which pin you need.
\end{warn}

\section{Using it}

\begin{lstlisting}
%%matlab -i df -o beta
fit = fitlm(df, 'y ~ x1 + x2');
beta = fit.Coefficients.Estimate;
disp(fit)
\end{lstlisting}

\texttt{-i} pushes a DataFrame in as a MATLAB \texttt{table}; \texttt{-o} brings
a variable back.

\begin{note}
\textbf{MATLAB is slow to start} --- roughly a minute cold. The first cell that
uses it pays for the whole session and every later one is fast. That is also why
\texttt{\%econ status} never starts it.

If MATLAB is already open, share it and EconEnv attaches instead of launching a
second copy: run \texttt{matlab.engine.shareEngine} in MATLAB, then
\texttt{\%econ config matlab.shared MATLAB\_shared}.
\end{note}

\section{Which release will start}

If you have several MATLAB releases installed, the one that starts is the one
your \emph{engine package} matches --- not the newest on disk.
\texttt{\%econ status} reports the release that will actually run.

% ======================================================================
\chapter{Exporting results for publication}

Getting a table out of a notebook and into a paper is where reproducibility
usually breaks: numbers are retyped, a coefficient is rounded twice, a
regression is re-run and the manuscript is not. EconEnv exports from the result
object itself, so the file and the estimation cannot drift apart.

\section{One call}

\begin{lstlisting}
econenv.export(cmp, "paper/table1", formats=["tex", "docx", "xlsx"])
\end{lstlisting}

writes \texttt{table1.tex}, \texttt{table1.docx} and \texttt{table1.xlsx} from
the same object. From a cell, importing nothing:

\begin{lstlisting}
%econ export cmp paper/table1 --formats tex,docx --caption "Table 1"
\end{lstlisting}

What you pass can be a single model, a comparison across five engines, a
DataFrame, a figure, or a list mixing them.

\section{Two layouts}

\subsection*{Journal --- the default}

What a paper prints: one column per model, the coefficient with significance
stars, the standard error beneath it in parentheses, and a foot of sample size
and fit statistics.

\begin{lstlisting}
econenv.export(cmp, "table1.tex", caption="Consumption function",
               label="tab:cons")
\end{lstlisting}

produces the following --- and this output was compiled with pdflatex during
development, which is how the escaping rules were settled:

\begin{center}
\begin{tabular}{@{}lccccc@{}}
\toprule
 & Python & EViews & MATLAB & R & Stata \\
\midrule
\_cons & 1.4805*** & 1.4805*** & 1.4805*** & 1.4805*** & 1.4805*** \\
 & (0.0350) & (0.0350) & (0.0350) & (0.0350) & (0.0350) \\
x1 & 1.9331*** & 1.9331*** & 1.9331*** & 1.9331*** & 1.9331*** \\
 & (0.0395) & (0.0395) & (0.0395) & (0.0395) & (0.0395) \\
Observations & 120 & 120 & 120 & 120 & 120 \\
R\textsuperscript{2} & 0.9597 & 0.9597 & 0.9597 & 0.9597 & 0.9597 \\
\bottomrule
\end{tabular}
\end{center}

\subsection*{Full --- every statistic}

One row per term per engine, with coefficient, standard error, test statistic,
p-value and confidence interval as separate columns. For your own checking
rather than for a paper.

\begin{lstlisting}
econenv.export(cmp, "check.xlsx", style="full")
\end{lstlisting}

\section{Formats}

\begin{center}
\begin{tabular}{@{}lllp{5.0cm}@{}}
\toprule
\textbf{Format} & \textbf{Ext} & \textbf{Needs} & \textbf{Notes} \\
\midrule
LaTeX & \texttt{.tex} & --- & booktabs rules, ready to input \\
Word & \texttt{.docx} & \texttt{export} extra & a real editable table, not a picture \\
Excel & \texttt{.xlsx} & \texttt{export} extra & one sheet per table, full precision \\
CSV & \texttt{.csv} & --- & \\
HTML & \texttt{.html} & --- & self-contained, styled like a journal table \\
Markdown & \texttt{.md} & \texttt{export} extra & \\
RTF & \texttt{.rtf} & --- & for submission systems that still want it \\
\bottomrule
\end{tabular}
\end{center}

\begin{lstlisting}
pip install "econenv[export]"
\end{lstlisting}

A filename with a suffix selects that format; a stem gives you the journal
bundle of \texttt{tex}, \texttt{docx} and \texttt{xlsx}.

\section{Significance stars}

The economics convention is the default --- \texttt{*** p<0.01},
\texttt{** p<0.05}, \texttt{* p<0.1} --- and the matching note is written under
the table automatically. For a journal that uses different thresholds:

\begin{lstlisting}
econenv.export(cmp, "t.tex",
               stars=((0.001, "***"), (0.01, "**"), (0.05, "*")))
\end{lstlisting}

\section{Figures, and why vector matters}

A PNG is fixed pixels: print it at column width and it looks fine, print it
larger and it does not. A PDF or EPS is instructions, so it stays sharp at any
size --- which is why most economics journals ask for vector figures.

Ask the engine for vector output \emph{before} running the cell, because a
bitmap cannot be converted into one afterwards:

\begin{lstlisting}
%econ config matlab.graphics pdf
%econ config r.graphics pdf
%econ config eviews.graphics svg
\end{lstlisting}

\begin{note}
EconEnv writes figures in whatever the engine actually produced, and the
extension follows the \emph{content} rather than the filename --- naming a PNG
\texttt{.pdf} makes a file nothing can open. If you export raster when vector
was wanted, the result says so and names the setting to change, rather than
wrapping a bitmap in a PDF container and calling it vector.
\end{note}

\section{What it will not do}

\begin{itemize}
  \item \textbf{Invent a statistic.} A number the engine did not report is left
        blank, not computed here under different assumptions. A table that
        quietly mixes two engines' conventions is worse than one with a gap.
  \item \textbf{Fake vector output.}
  \item \textbf{Round twice.} The journal layout formats once from the
        full-precision value; the Excel export keeps the unrounded number.
\end{itemize}

% ======================================================================
\chapter{Moving data between programs}

This is the feature the whole project exists for: build a dataset once and use
it everywhere, without a single CSV.

\section{One dataset, every engine}

\begin{lstlisting}
econenv.broadcast("macro", df)
\end{lstlisting}

\begin{lstlisting}
{'eviews': None, 'matlab': None, 'r': None, 'stata': None}
\end{lstlisting}

\texttt{None} means it arrived. An engine that is missing or fails is recorded
rather than stopping the rest, because a machine without MATLAB should still get
the data into R and Stata. Pass \texttt{strict=True} to raise instead, or name
the engines you want.

\section{One at a time}

\begin{lstlisting}
econenv.push("stata", "macro", df)          # Python to Stata
back = econenv.pull("r", "macro")           # R to Python
econenv.move("stata", "matlab", "macro")    # Stata to MATLAB
\end{lstlisting}

Or inline, while running code anyway:

\begin{lstlisting}
%%R -i macro -o results
%%eviews -i macro
%%matlab -i macro -o beta
\end{lstlisting}

\section{What can move where}

\begin{lstlisting}
econenv.transfer_matrix()
\end{lstlisting}

reports what is possible on \emph{this} machine rather than in principle. An
engine that is installed but not configured cannot take your data, and learning
that before you start is cheaper than a failure mid-session.

\section{What survives the trip}

\begin{center}
\begin{tabular}{@{}lp{9.2cm}@{}}
\toprule
\textbf{Type} & \textbf{Behaviour} \\
\midrule
Numeric & exact \\
Dates & preserved; the index becomes a column in R and MATLAB, and a dated page in EViews \\
Booleans & preserved, except that MATLAB logicals cannot be missing \\
Categoricals & become factors in R, categoricals in MATLAB \\
Strings & preserved; missing strings become empty in MATLAB \\
Missing values & preserved, with Stata's extended missing values noted \\
\bottomrule
\end{tabular}
\end{center}

Anything lossy warns, naming the column and what changed. Anything that would
alter a value raises instead.

\begin{warn}
\textbf{EViews names work differently.} Pushing to EViews creates a \emph{page
of series}, not an object named after the frame --- there is no \texttt{macro}
object in the workfile, only \texttt{X}, \texttt{Y}, \texttt{Z}. EconEnv
remembers the name you pushed under so \texttt{pull("eviews", "macro")} returns
the page, but inside a \texttt{\%\%eviews} cell you refer to the series by their
own names.
\end{warn}

% ======================================================================
\chapter{Google Colab}

Colab is Linux, and that settles what is possible before anything is installed.

\section{What runs on a Colab cloud runtime}

\begin{center}
\begin{tabular}{@{}llp{7.4cm}@{}}
\toprule
\textbf{Engine} & \textbf{Colab} & \textbf{Why} \\
\midrule
Python & works & it is the kernel \\
R & works & R is already on the Colab image \\
Stata & \textbf{possible} & Stata for Linux exists and pystata supports it \\
EViews & no & no Linux build, and no permitted workaround \\
\bottomrule
\end{tabular}
\end{center}

Install as usual:

\begin{lstlisting}
%pip install econenv
\end{lstlisting}

\section{Stata on a Colab runtime}

This does work, if you hold a Stata \textbf{Linux} licence. Keep the Linux
tarball and your \texttt{stata.lic} in Google Drive, mount Drive, and install at
the top of the notebook:

\begin{lstlisting}
from google.colab import drive
drive.mount('/content/drive')
\end{lstlisting}

\begin{lstlisting}
!mkdir -p /usr/local/stata19
!tar -xzf "/content/drive/MyDrive/stata/Stata19Linux64.tar.gz" -C /usr/local/stata19
!cd /usr/local/stata19 && yes | ./install
!cp "/content/drive/MyDrive/stata/stata.lic" /usr/local/stata19/
\end{lstlisting}

EconEnv finds it with no configuration: the Linux executable names
(\texttt{stata-mp}, \texttt{stata-se}, \texttt{stata}) and \texttt{/usr/local}
have been part of discovery from the start.

\begin{warn}
Three caveats, none of them technical. You need a \textbf{Linux} licence
entitlement --- a Windows-only licence does not cover this. The runtime is
destroyed at the end of the session, so the install repeats every time. And
whether your licence permits installation on a disposable cloud VM is a question
for StataCorp, not for this project.
\end{warn}

\section{EViews on a Colab runtime --- not possible}

Three independent blockers, none of which EconEnv can route around:

\begin{enumerate}
  \item \textbf{There is no Linux build.} EViews ships for Windows and macOS.
  \item \textbf{Wine does not work.} EViews cannot read a valid machine ID under
        Wine, so licence activation fails.
  \item \textbf{Remote access is forbidden.} The obvious workaround --- run
        EViews on your Windows machine and reach it from Colab over a tunnel ---
        is ruled out by EViews' own documentation, which states that
        \emph{``web server access to EViews via COM is not allowed''} and limits
        remote Distributed COM to a single instance.
\end{enumerate}

The third is the one that settles it: that bridge is straightforward to build
and contractually prohibited, so EconEnv does not ship it.

\section{All four engines, still in Colab}

There is a way to have the Colab interface \emph{and} every engine, and it
avoids all of the above.

Colab runs in a browser \textbf{on your own PC}. Google lets you point that
interface at a Jupyter server running on that same PC --- a \textbf{local
runtime}. The notebook UI stays Colab; the kernel, and therefore Python, R,
Stata and EViews, is your Windows machine.

\begin{note}
Nothing is exposed to the internet. Your browser talks to \texttt{localhost},
Google's servers never reach your machine, and EViews is driven by local COM
exactly as it would be in a local notebook. There is no web server in front of
EViews, so this is not the restriction quoted above.
\end{note}

It needs the \textbf{classic} Jupyter stack. The bridge package
\texttt{jupyter\_http\_over\_ws} was last released in March 2020 and is a
notebook 5/6 server extension; on \texttt{notebook 7} and
\texttt{jupyter\_server 2} it does not load at all.

\begin{center}
\begin{tabular}{@{}lll@{}}
\toprule
\textbf{Stack} & \textbf{Enable step} & \texttt{/http\_over\_websocket} \\
\midrule
notebook 7.5.5 & fails & 404 \\
notebook 6.5.7 & fails & 404 \\
\textbf{notebook 6.4.12} & \textbf{validates} & \textbf{HTTP 400} \\
\bottomrule
\end{tabular}
\end{center}

HTTP 400 is the correct answer: the endpoint exists and is waiting for the
websocket upgrade Colab performs. Both other stacks run on
\texttt{jupyter\_server 2}, which will not load a legacy server extension.

So use a dedicated environment:

\begin{lstlisting}
python -m venv colab-runtime
colab-runtime/Scripts/pip install notebook==6.4.12 jupyter_http_over_ws econenv
colab-runtime/Scripts/jupyter serverextension enable --py jupyter_http_over_ws
colab-runtime/Scripts/jupyter notebook --no-browser --port=8888
    --NotebookApp.port_retries=0
    --NotebookApp.allow_origin="https://colab.research.google.com"
\end{lstlisting}

Copy the \texttt{http://localhost:8888/?token=...} line it prints. In Colab,
click the \textbf{Connect} arrow, choose \textbf{Connect to a local runtime},
and paste the URL.

Then everything works, because the kernel is your PC:

\begin{lstlisting}
%load_ext econenv
%econ status
\end{lstlisting}

Two practical caveats: your PC must stay awake with the server running for as
long as the notebook is open, and that environment is separate from your usual
one, so install into it whatever the notebook needs.

% ======================================================================
\chapter{A worked example}

The notebook \texttt{examples/10\_real\_data\_four\_engines.ipynb} works through a
complete exercise on \textbf{real US quarterly macroeconomic data, 1959Q1 to
2009Q3} --- 203 observations, shipped with statsmodels, so it needs no download.

\section{What it does}

\begin{center}
\begin{tabular}{@{}llp{6.2cm}@{}}
\toprule
\textbf{Step} & \textbf{Engine} & \textbf{Why that one} \\
\midrule
Load and transform & Python & pandas is the best tool for this \\
Describe, plot, diagnose & R & \texttt{plot(fit)} and R's diagnostic vocabulary \\
HAC standard errors & Stata & \texttt{newey} is the reference implementation \\
Unit roots, cointegration & EViews & this is what EViews is for \\
Forecast and stability & EViews & CUSUM, forecast bands \\
Compare all four & EconEnv & one specification, four engines, one table \\
\bottomrule
\end{tabular}
\end{center}

\section{The model}

A textbook consumption function,
\[
  \log C_t = \beta_0 + \beta_1 \log Y_t + \beta_2 r_t + \varepsilon_t ,
\]
with $C$ real consumption, $Y$ real GDP and $r$ the real interest rate.

Both series trend strongly, so the interesting questions --- are they
non-stationary, are they cointegrated, are the standard errors trustworthy ---
are exactly the ones each program answers differently. That is what makes it a
good test of the whole idea.

\section{A representative result}

EViews, on log real GDP:

\begin{lstlisting}
Null Hypothesis: LRGDP has a unit root
Exogenous: Constant
Lag Length: 1 (Automatic - based on SIC, maxlag=14)

                                        t-Statistic   Prob.*
Augmented Dickey-Fuller test statistic   -1.820451    0.3698
Test critical values:    1% level        -3.462901
                         5% level        -2.875752
\end{lstlisting}

and the four-engine comparison of the same regression:

\begin{lstlisting}
Coefficients
              python         eviews              r          stata
term
_cons      1.4805338      1.4805338      1.4805338      1.4805338
lrgdp      0.8863268      0.8863268      0.8863268      0.8863268
realint   -0.0936228     -0.0936228     -0.0936228     -0.0936228
\end{lstlisting}

Maximum disagreement across the four: of the order of $10^{-15}$ --- floating
point noise and nothing more.

\section{What deliberately does not agree}

The information criteria differ, and EconEnv reports why rather than choosing
one:

\begin{itemize}
  \item \texttt{statsmodels}: $-2\ell + 2k$
  \item R: counts $\sigma^2$ as a parameter, so $k+1$
  \item EViews: divides by $n$
  \item Stata: from \texttt{estat ic}, with $k = e(\text{rank}) + 1$
\end{itemize}

None is wrong. They answer slightly different questions, and a comparison that
hid that would be worse than useless.

% ======================================================================
\chapter{Reproducibility}

\section{Snapshots}

\begin{lstlisting}
econenv.snapshot()
\end{lstlisting}

records the EconEnv version, the Python version and packages, the R version,
the Stata version and edition, the EViews version, the operating system,
architecture and timestamp. Export it as JSON and keep it with your results.

\section{Why this matters here more than usual}

A result produced across four programs has four times as many ways to become
irreproducible. A snapshot taken at the time of estimation is the difference
between ``we used Stata'' and ``we used StataNow 19.5 MP with PyStata, on
Windows, in March''.

% ======================================================================
\chapter{Troubleshooting}

\section{The extension will not load}
\label{sec:wrongpython}

\begin{lstlisting}
ModuleNotFoundError: No module named 'econenv'
\end{lstlisting}

Almost always: Jupyter is running on a different Python from the one you
installed EconEnv into. Check from inside the notebook:

\begin{lstlisting}
import sys; print(sys.executable)
\end{lstlisting}

then install into exactly that interpreter:

\begin{lstlisting}
!{sys.executable} -m pip install econenv
\end{lstlisting}

\section{An engine will not start}

\begin{lstlisting}
%econ doctor
\end{lstlisting}

Every failing check names a fix. The usual causes:

\begin{itemize}
  \item \textbf{EViews}: never opened by hand, so its COM server is not
        registered. Open it once and close it.
  \item \textbf{EViews}: \texttt{comtypes} not installed ---
        \texttt{pip install "econenv[eviews]"}.
  \item \textbf{Stata}: version 16 or earlier. PyStata requires Stata~17.
  \item \textbf{R}: installed somewhere unusual --- set \texttt{r.home}.
\end{itemize}

\section{EViews reports the wrong version}

With several versions installed, the generic identifier binds to whichever
registered last. \texttt{\%econ status} always shows the one actually connected.
Pin a version with \texttt{\%econ config eviews.progid EViews.Manager.14}.

\section{rpy2 is installed but broken}
\label{sec:rpy2}

A typical symptom:

\begin{lstlisting}
ImportError: cannot import name 'SexpVectorCCompatibleAbstract'
    from 'rpy2.rinterface_lib.sexp'
\end{lstlisting}

This means a version built for an older Python. PyPI has no Windows wheels for
rpy2, so \texttt{pip install rpy2} cannot fix it. Use conda-forge:

\begin{lstlisting}
pip uninstall -y rpy2
conda install -c conda-forge rpy2
\end{lstlisting}

\begin{warn}
conda-forge's rpy2 brings \emph{its own R} with it, separate from any R already
installed. Packages in your existing R will not be visible to it. Unless you
specifically need rpy2, stay on EconEnv's subprocess backend, which works
without any of this.
\end{warn}

\section{A cell runs and shows nothing}

Report it. Every display command should produce output, a plot, or a warning.
Silence is a bug:
\url{https://github.com/merwanroudane/econenv/issues}

% ======================================================================
\chapter{Licensing and citation}

\section{EconEnv}

MIT licence. Free for any use, including commercial.

\section{The commercial programs}

EconEnv's licence says nothing about Stata or EViews. It contains no part of
either, and connects only to installations you have licensed yourself. Your
obligations under those licences are unchanged.

\section{Developer}

\textbf{Dr Merwan Roudane} --- author and maintainer of EconEnv.

\begin{itemize}
  \item GitHub: \url{https://github.com/merwanroudane}
  \item Repository: \url{https://github.com/merwanroudane/econenv}
  \item Package: \url{https://pypi.org/project/econenv/}
  \item Issues and feature requests: \url{https://github.com/merwanroudane/econenv/issues}
\end{itemize}

Contributions are welcome. If EconEnv reports something incorrectly, or a
command runs and shows nothing, that is a bug worth reporting --- silence is
never the intended behaviour.

\section{Citation}

\begin{lstlisting}
Roudane, M. (2026). EconEnv: One Notebook, Multiple Econometric
Engines. Python package. https://github.com/merwanroudane/econenv
\end{lstlisting}

A \texttt{CITATION.cff} file in the repository is kept up to date, and GitHub
will format a citation from it automatically.

\vfill
\begin{center}
\rule{0.4\textwidth}{0.4pt}\\[8pt]
\textbf{EconEnv} --- Dr Merwan Roudane\\
\url{https://github.com/merwanroudane/econenv}
\end{center}

\end{document}
